\documentclass[11pt]{amsart}
\usepackage[T1]{fontenc}
\usepackage{amsmath,amssymb,amsthm}
\usepackage[margin=1.1in]{geometry}
\usepackage{booktabs}
\usepackage{enumitem}
\usepackage{hyperref}

\newtheorem{theorem}{Theorem}[section]
\newtheorem{lemma}[theorem]{Lemma}
\newtheorem{proposition}[theorem]{Proposition}
\newtheorem{corollary}[theorem]{Corollary}
\newtheorem{conjecture}[theorem]{Conjecture}
\theoremstyle{definition}
\newtheorem{definition}[theorem]{Definition}
\newtheorem{remark}[theorem]{Remark}
\newtheorem{example}[theorem]{Example}

\title{A direct lift to the two-end test through triangles at the even vertex}
\author{Gallai project}
\date{27 November 2026}

\begin{document}

\begin{abstract}
The two-end test (T1) of the spine reduction asks, for a graph $J$ of odd order $2k+1$ whose only
even vertex is $z$ and for a neighbour $r$ of $z$, for a decomposition of $J-z$ into at most $k$ paths
with at least two ends at $r$. We show that (T1) follows from the existence of a single Lov\'asz
decomposition of an auxiliary graph $M$ with one even vertex, subject to one \emph{disjunctive}
end-path condition per triangle $z a_i b_i$. Here $M$ is obtained from $J$ by deleting an edge $zw$
and the triangles $za_ib_i$ of a near-perfect matching of $N(z)-r$. The lift does not pass through
star-terminal decompositions. For $d(z)=6$ it replaces the five or six end-path conditions of the
dense-residual base conjecture by two disjunctive ones. Exact computation supports the resulting base
conjecture on random dense $3$-connected graphs. The reduction was proposed by an external
collaborator; the proof and data here are an independent check.
\end{abstract}

\maketitle

\section{Setting}

All graphs are finite and simple. A \emph{Lov\'asz decomposition} of a graph $H$ with at most one
vertex of positive even degree is a partition of $E(H)$ into paths of positive length such that every
vertex $v$ is an end of exactly $d_H(v) \bmod 2$ of them. Lov\'asz~\cite{L68} proved that one always
exists. If $v$ has odd degree, $Q(v)$ denotes the unique path ending at $v$.

Throughout, $J$ has order $2k+1$, minimum degree at least $3$, and exactly one vertex $z$ of even
degree, with $d_J(z)=2m\ge 4$. The test (T1) at $(J,z,r)$, for $r\in N_J(z)$, asks for a decomposition
of $B:=J-z$ into at most $k$ paths with at least two ends at $r$ (Lean: \texttt{Gallai.Test1}).

\begin{definition}[triangle data]\label{def:data}
Fix $r\in N_J(z)$, a vertex $w\in N_J(z)-r$, and a perfect matching
$M_0=\{a_1b_1,\dots,a_{m-1}b_{m-1}\}$ of $N_J(z)-\{r,w\}$ consisting of edges of $J$. Put
\[
M \;=\; J-\bigl(\{zw\}\cup\{za_i,\ zb_i,\ a_ib_i : 1\le i\le m-1\}\bigr),
\]
deleting $3m-2$ edges.
\end{definition}

\begin{lemma}\label{lem:M}
In $M$, the vertex $z$ is a leaf adjacent to $r$; the vertex $w$ is the unique vertex of positive even
degree, and $d_M(w)\ge 2$; every other vertex has odd degree. Hence every Lov\'asz decomposition of $M$
has exactly $k$ paths.
\end{lemma}

\begin{proof}
Among the $2m$ edges at $z$ only $zr$ survives. The vertex $w$ loses one edge, so it has even degree
$d_J(w)-1\ge 2$. Each $a_i$ and $b_i$ loses two edges ($za_i,a_ib_i$, resp.\ $zb_i,a_ib_i$), so its
degree stays odd, and it is at least $1$. All other degrees are unchanged, and $z$ has degree $1$. So
the $2k$ vertices other than $w$ are odd. Each is an end of exactly one path, so there are $k$ paths.
Disconnected $M$ is handled componentwise.
\end{proof}

\section{The lift}

\begin{theorem}[direct lift]\label{thm:lift}
With the data of Definition~\ref{def:data}, suppose $M$ has a Lov\'asz decomposition $D_0$ such that,
for every $i$,
\begin{equation}\label{eq:disj}
a_i\notin V(Q_0(b_i)) \quad\text{or}\quad b_i\notin V(Q_0(a_i)).
\end{equation}
Then $(J,z,r)$ passes (T1).
\end{theorem}

\begin{proof}
\emph{Step 1: delete $z$.} Let $B_0=B-M_0=M-z$. In $B_0$ the vertex $w$ has even degree, every $a_i$,
$b_i$ has odd degree, and $r$ has even degree $d_J(r)-1\ge 2$. The path $Q_0(z)$ begins $z,r$.
\begin{itemize}[leftmargin=*]
\item If $Q_0(z)$ has length at least $2$, replace it by $Q_0(z)-z$, which starts at $r$. The vertex
$r$ is interior to $Q_0(z)$, so its unique end in $D_0$ lies on another path $Q_0(r)$. Now $r$ is an end
of two paths.
\item If $Q_0(z)=zr$, delete it. Then $r$ is an end of no path, and since $r$ has positive even degree
in $B_0$, some path passes through $r$ internally. Split that path at $r$ into two paths of positive
length, both ending at $r$.
\end{itemize}
In both cases we obtain a decomposition $D_1$ of $B_0$ into $k$ paths of positive length. In $D_1$, $r$
is an end of exactly two paths. Every odd vertex $v$ of $B_0$ other than $r$ is an end of exactly one
path $P_1(v)$, and $V(P_1(v))\subseteq V(Q_0(v))$, because each step only deletes $z$ or splits a path.

\emph{Step 2: restore the matching.} By~\eqref{eq:disj} and the inclusion just proved, for each $i$ we
may choose $(v_i,u_i)\in\{(b_i,a_i),(a_i,b_i)\}$ with $u_i\notin V(P_1(v_i))$. Process $i=1,\dots,m-1$ in
turn: append the edge $v_iu_i$ to the end $v_i$ of the current path $P(v_i)$. This is legal because
$u_i\notin V(P(v_i))$. Indeed, $V(P(v_i))$ equals $V(P_1(v_i))$ unless the same path was earlier
extended at its other end by some $u_j$, and $u_j\ne u_i$ because $M_0$ is a matching. A path has two
ends, so it is extended at most twice. The number of paths stays $k$, and every edge of $B=B_0+M_0$ is
used once. No $a_i$, $b_i$ equals $r$, so $r$ still ends two paths. This is (T1) at $(J,z,r)$.
\end{proof}

\begin{remark}[comparison]\label{rem:compare}
For $m=3$ the dense-residual base conjecture of the T(6)-i manuscript
(\texttt{manuscripts/t6-dense}, Conj.~5.1; Lean \texttt{test1\_of\_appendBases}) needs five or six
end-path conditions on its base. Audit A26 shows its main branch is equivalent to the star-terminal
condition STL-P plus two prescribed terminal edges, a strengthening of STL. Theorem~\ref{thm:lift}
bypasses star-terminal decompositions and needs only two disjunctive conditions. Neither base condition
obviously implies the other, since the auxiliary graphs differ. Both are sufficient for (T1).
\end{remark}

\begin{remark}[availability of the data]
Definition~\ref{def:data} needs a perfect matching of $J[N(z)-\{r,w\}]$ for some $w$. In the dense
regime of T(6)-i, where $J[N(z)]=K_6-D$ with $D\in\{\emptyset,K_2,P_3,K_3,C_4\}$, such a matching exists
for every root. If $J[N(z)-r]$ has no near-perfect matching, one can still peel the triangles of a
maximal matching. Peeling leaves an instance in which $N(z)-r$ is independent, where pairings of
$N(z)-r$ use non-edges only.
\end{remark}

\begin{conjecture}[triangle base]\label{conj:base}
If $J$ is $3$-connected, then for every $r$ there are data as in Definition~\ref{def:data} and a
Lov\'asz decomposition of $M$ satisfying~\eqref{eq:disj} for all $i$.
\end{conjecture}

By Theorem~\ref{thm:lift}, Conjecture~\ref{conj:base} implies (T1) for every $3$-connected $J$ whose
even vertex admits the data. With the spine reduction, this covers the two-end half of Gallai's
conjecture in that range.

\section{A local guard}

The disjunction~\eqref{eq:disj} fails at a pair $\{a,b\}$ exactly when $Q_0(a)=Q_0(b)$, or when each of
$Q_0(a),Q_0(b)$ contains the other vertex. The next example shows that a repair cannot work with these
two paths alone.

\begin{example}\label{ex:guard}
Let $H$ be the $4$-cycle $a\,x\,b\,y$ with pendant edges $bt_1$ and $at_2$, so that $ab\notin E(H)$. The
odd vertices are $a,b,t_1,t_2$ and $x,y$ have degree $2$, so a decomposition with one end at each odd
vertex has two paths. Since $t_1$ is a leaf at $b$, no path joins $t_1$ to $b$; similarly for $t_2$ and
$a$. So the ends are paired either as $\{t_1,t_2\},\{a,b\}$ or as $\{t_1,a\},\{t_2,b\}$.
\begin{itemize}[leftmargin=*]
\item In the first case the paths are $t_1\,b\,x\,a\,t_2$ and $b\,y\,a$ (or with $x,y$ exchanged), and
$Q(a)=Q(b)=b\,y\,a$.
\item In the second case they are $t_1\,b\,x\,a$ and $t_2\,a\,y\,b$ (or with $x,y$ exchanged), and each
of $Q(a),Q(b)$ contains the other vertex.
\end{itemize}
Thus every such decomposition violates~\eqref{eq:disj}. A local uncrossing inside $M$ must therefore
involve a third path, ending at an odd vertex.
\end{example}

\section{Evidence}

The script \texttt{scripts/dtl.py} generates random $J$ with $d(z)=6$, $J[N(z)]$ dense (each pair an
edge with probability $0.9$ before parity repair), minimum degree $\ge 3$ and connectivity $\ge 3$.
For every root $r$ it searches for data and a base decomposition with the exact oracle
\texttt{tools/route\_ext.c}: the four orientation choices of~\eqref{eq:disj} are four per-end avoidance
queries, with a node budget per query. On success it executes Steps~1--2 and verifies the result
independently. The check confirms an edge partition of $J-z$ into simple paths, at most $k$ paths,
and at least two ends at $r$.

\begin{center}
\begin{tabular}{lrrrrr}
\toprule
order, $d(z)$ & graphs & roots & base found & lift verified & first data choice works\\
\midrule
$11$, $6$ & $30$ & $180$ & $180$ & $180$ & $175$\\
\bottomrule
\end{tabular}
\end{center}

Nine individual queries exceeded the budget; in each case another data choice at the same root
succeeded. No lift failed. The generalised script \texttt{scripts/dtl2.py} handles $d(z)=2m$ with
$m-1$ pairs; results at $d(z)=8$ are recorded in \texttt{scripts/out/}.

\bibliographystyle{plain}
\begin{thebibliography}{9}
\bibitem{L68} L.~Lov\'asz, On covering of graphs, in: \emph{Theory of Graphs} (Proc.\ Colloq.,
Tihany, 1966), Academic Press, New York, 1968, pp.~231--236.
\end{thebibliography}

\end{document}
